\clearpage
\subsection{Page-Table Entry Format}\label{section:page-table-entry-format}

Every page-table entry is one naturally aligned 64-bit descriptor. Bits
\texttt{55..14} are a fixed (:term:base.terms.physical_frame_number:) (:term:base.terms.field:), so this format has a
56-bit architectural physical-address ceiling. An implementation reports its
exact (:ref:base.cpuid.IMPLEMENTATION.ADDRESS_WIDTHS.PARAMETERS.PABITS:) value, from 32 through 56, through CPUID
\texttt{ADDRESS\_WIDTHS}. In a present descriptor, (:ref:base.memory.translation.PTE.PFN:) bits at or above that
value (:term:base.terms.must_be_zero:). Bits \texttt{13..11} are (:term:base.terms.reserved:) and (:term:base.terms.must_be_zero:) in a present descriptor.
Bits \texttt{57..56} are (:ref:base.memory.translation.PTE.SW:)\texttt{[1:0]} and are (:term:base.terms.ignored:)
by hardware; bits \texttt{63..58} are (:term:base.terms.reserved:) and (:term:base.terms.must_be_zero:).

(:pte-field-target:base.memory.translation.PTE.P:)\texttt{P} in bit 0 is common. When (:ref:base.memory.translation.PTE.P:)\texttt{=0}, the remaining bits are
(:term:base.terms.ignored:) and the entry is not present. When (:ref:base.memory.translation.PTE.P:)\texttt{=1}, (:pte-field-target:base.memory.translation.PTE.T:)\texttt{T} in bit 1
selects the leaf or table-pointer layout. A nonzero (:term:base.terms.reserved:) bit in a present
descriptor raises (:ref:base.events.EXCEPTION.MALFORMED_PAGE_TABLE_ENTRY:).

\BedrockTableCaption{Leaf Descriptor Fields ((:ref:base.memory.translation.PTE.P:)\texttt{=1}, (:ref:base.memory.translation.PTE.T:)\texttt{=0})}
\begin{BedrockTabular}{@{}>{\raggedright\arraybackslash}p{0.75in}>{\raggedright\arraybackslash}p{0.75in}>{\raggedright\arraybackslash}p{4.00in}@{}}
\toprule
\rowcolor{BedrockHeaderFill}\textbf{Field} & \textbf{Bits} & \textbf{Meaning}\\
\midrule
(:ref:base.memory.translation.PTE.P:) & \texttt{0} & present, one\\
(:ref:base.memory.translation.PTE.T:) & \texttt{1} & leaf selector, zero\\
(:pte-field-target:base.memory.translation.PTE.AM:)\texttt{AM} & \texttt{4..2} & permission and Normal/MMIO access class\\
(:pte-field-target:base.memory.translation.PTE.U:)\texttt{U} & \texttt{5} & user-domain permission\\
(:pte-field-target:base.memory.translation.PTE.G:)\texttt{G} & \texttt{6} & global translation\\
(:pte-field-target:base.memory.translation.PTE.A:)\texttt{A} & \texttt{7} & accessed state\\
(:pte-field-target:base.memory.translation.PTE.D:)\texttt{D} & \texttt{8} & exact dirty state\\
(:pte-field-target:base.memory.translation.PTE.CP:)\texttt{CP} & \texttt{10..9} & cache policy\\
(:term:base.terms.reserved:) & \texttt{13..11} & (:term:base.terms.must_be_zero:)\\
(:pte-field-target:base.memory.translation.PTE.PFN:)\texttt{PFN} & \texttt{55..14} & mapped (:term:base.terms.physical_frame_number:)\\
(:pte-field-target:base.memory.translation.PTE.SW:)\texttt{SW[1:0]} & \texttt{57..56} & software-defined, (:term:base.terms.ignored:) by hardware\\
(:term:base.terms.reserved:) & \texttt{63..58} & (:term:base.terms.must_be_zero:)\\
\bottomrule
\end{BedrockTabular}

\BedrockTableCaption{Leaf (:ref:base.memory.translation.PTE.AM:) Encodings}
\begin{BedrockTabular}{@{}>{\raggedright\arraybackslash}p{0.65in}>{\raggedright\arraybackslash}p{1.35in}>{\raggedright\arraybackslash}p{1.15in}>{\raggedright\arraybackslash}p{2.35in}@{}}
\toprule
\rowcolor{BedrockHeaderFill}\textbf{(:ref:base.memory.translation.PTE.AM:)} & \textbf{Permission} & \textbf{Class} & \textbf{Name}\\
\midrule
\texttt{000} & read & Normal & \texttt{R}\\
\texttt{001} & write & Normal & \texttt{W}\\
\texttt{010} & execute & Normal & \texttt{X}\\
\texttt{011} & read/write & Normal & \texttt{RW}\\
\texttt{100} & read/execute & Normal & \texttt{RX}\\
\texttt{101} & read & MMIO & \texttt{MMIO\_R}\\
\texttt{110} & write & MMIO & \texttt{MMIO\_W}\\
\texttt{111} & read/write & MMIO & \texttt{MMIO\_RW}\\
\bottomrule
\end{BedrockTabular}

No leaf encoding grants both write and execute, and no MMIO leaf grants
execute. This W\textasciicircum X property applies to one selected translation.
The architecture does not compare aliases: distinct descriptors for the same
physical page may have different permissions, (:ref:base.memory.translation.PTE.AM:) classes, or (:ref:base.memory.translation.PTE.CP:) values.

\BedrockTableCaption{Table-Pointer Descriptor Fields ((:ref:base.memory.translation.PTE.P:)\texttt{=1}, (:ref:base.memory.translation.PTE.T:)\texttt{=1})}
\begin{BedrockTabular}{@{}>{\raggedright\arraybackslash}p{0.75in}>{\raggedright\arraybackslash}p{0.75in}>{\raggedright\arraybackslash}p{4.00in}@{}}
\toprule
\rowcolor{BedrockHeaderFill}\textbf{Field} & \textbf{Bits} & \textbf{Meaning}\\
\midrule
(:ref:base.memory.translation.PTE.P:) & \texttt{0} & present, one\\
(:ref:base.memory.translation.PTE.T:) & \texttt{1} & table-pointer selector, one\\
(:pte-field-target:base.memory.translation.PTE.R:)(:pte-field-target:base.memory.translation.PTE.W:)(:pte-field-target:base.memory.translation.PTE.X:)\texttt{R/W/X} & \texttt{2/3/4} & subtree maximum permissions\\
(:ref:base.memory.translation.PTE.U:) & \texttt{5} & subtree maximum user-domain permission\\
(:term:base.terms.reserved:) & \texttt{6} & (:term:base.terms.must_be_zero:)\\
(:ref:base.memory.translation.PTE.A:) & \texttt{7} & accessed state\\
(:term:base.terms.reserved:) & \texttt{13..8} & (:term:base.terms.must_be_zero:)\\
(:ref:base.memory.translation.PTE.PFN:) & \texttt{55..14} & next-table (:term:base.terms.physical_frame_number:)\\
(:ref:base.memory.translation.PTE.SW:)\texttt{[1:0]} & \texttt{57..56} & software-defined, (:term:base.terms.ignored:) by hardware\\
(:term:base.terms.reserved:) & \texttt{63..58} & (:term:base.terms.must_be_zero:)\\
\bottomrule
\end{BedrockTabular}

A table pointer with (:ref:base.memory.translation.PTE.R:)\texttt{=}(:ref:base.memory.translation.PTE.W:)\texttt{=}(:ref:base.memory.translation.PTE.X:)\texttt{=0}, or a table pointer at L1, is malformed.
The walker begins with (:ref:base.memory.translation.PTE.R:), (:ref:base.memory.translation.PTE.W:), (:ref:base.memory.translation.PTE.X:), and (:ref:base.memory.translation.PTE.U:) enabled, intersects each table bitmap and
(:ref:base.memory.translation.PTE.U:) bit, then intersects the terminating leaf\textquotesingle{}s decoded (:ref:base.memory.translation.PTE.AM:) permissions and (:ref:base.memory.translation.PTE.U:) bit.
The leaf alone supplies (:ref:base.memory.translation.PTE.G:), (:ref:base.memory.translation.PTE.CP:), and Normal/MMIO class. An L1, L2, L3, or L4 leaf must
have a physical base aligned to $2^{14}$, $2^{23}$, $2^{34}$, or $2^{45}$ bytes respectively.

The (:ref:base.memory.translation.PTE.CP:) (:term:base.terms.field:) is orthogonal to (:ref:base.memory.translation.PTE.AM:): all 32 (:ref:base.memory.translation.PTE.AM:)/(:ref:base.memory.translation.PTE.CP:) combinations are structurally
valid. (:ref:base.memory.translation.PTE.CP:) values 0 through 3 mean cacheable write-back, coherent uncacheable,
coherent write-through, and coherent write-combining respectively. (:ref:base.memory.translation.PTE.CP:) never
weakens an MMIO execution rule.

Hardware sets (:ref:base.memory.translation.PTE.A:) on each structurally valid descriptor consumed by a completed
walk and may also set (:ref:base.memory.translation.PTE.A:) during a speculative walk whose instruction later
faults or is squashed. (:ref:base.memory.translation.PTE.D:) is set only when the corresponding store or atomic update
commits, and every required (:ref:base.memory.translation.PTE.D:) update succeeds before its memory update becomes
visible. Hardware (:ref:base.memory.translation.PTE.A:)/(:ref:base.memory.translation.PTE.D:) changes are conditional atomic updates of the complete
64-bit PTE and preserve every other (:term:base.terms.field:). A concurrent change causes a
re-read and complete revalidation; stale (:term:base.terms.field|plural:) are never written back.
(:ref:base.instructions.PTQUERY:) and (:ref:base.instructions.VTOP:) do not modify (:ref:base.memory.translation.PTE.A:) or (:ref:base.memory.translation.PTE.D:).
